\chapter{Memory Systems}
\label{ch:memory}

\epigraph{``Memory is what turns a conversation into a relationship. For agents, it turns a session into a capability.''}{}

\learninggoal{
	\item Distinguish the four types of agent memory: working, episodic, semantic, and procedural.
	\item Implement context window management strategies for bounded-length conversations.
	\item Design a long-term memory system using vector storage and retrieval.
	\item Understand the trade-offs between summarisation, retrieval, and structured memory.
}

\section{A Taxonomy of Agent Memory}

Agents, like humans, need multiple memory systems operating at different timescales and levels of abstraction.

\begin{definition}[Memory Types]
	\begin{itemize}
		\item \textbf{Working memory:} the current context window. Fast, volatile, bounded ($\sim 10^5$ tokens).
		\item \textbf{Episodic memory:} records of past episodes (trajectories, successes, failures). Stored in vector databases or logs.
		\item \textbf{Semantic memory:} factual knowledge about the world and the user. Retrieved from vector stores or knowledge graphs.
		\item \textbf{Procedural memory:} knowledge of \emph{how to do things}---tool usage patterns, reasoning strategies. Stored as prompts or fine-tuned behaviours.
	\end{itemize}
\end{definition}

\begin{longtable}{p{3cm}p{3cm}p{3cm}p{3cm}p{2.5cm}}
	\toprule
	\textbf{Property} & \textbf{Working} & \textbf{Episodic} & \textbf{Semantic} & \textbf{Procedural} \\
	\midrule
	Timescale         & Seconds          & Hours--days       & Persistent        & Persistent          \\
	Capacity          & 100K tokens      & Unlimited         & Unlimited         & Unlimited           \\
	Retrieval         & Direct           & Similarity search & Similarity search & Prompt engineering  \\
	Volatility        & Session-only     & Archival          & Archival          & Updateable          \\
	Update            & Append-only      & Summarise + store & Add/merge facts   & Versioned prompts   \\
	\bottomrule
\end{longtable}

\section{Working Memory Management}

The context window is the agent's working memory. It is the highest-bandwidth memory system (all tokens attend to all others) but strictly bounded.

\subsection{The Context Window Budget}

Each token in the context window costs $O(d)$ in computation and $O(1)$ in memory (KV cache). The budget must be allocated across:

\begin{itemize}
	\item \textbf{System prompt:} instructions, tool definitions (typically 1K--5K tokens).
	\item \textbf{Message history:} conversation turns (grows unboundedly).
	\item \textbf{Tool outputs:} results of tool calls (highly variable, can be 10K+ tokens).
	\item \textbf{Retrieved context:} documents from long-term memory (variable).
\end{itemize}

\subsection{Eviction Policies}

When the context window is full, something must be removed:

\begin{longtable}{p{3.5cm}p{4cm}p{4cm}p{3cm}}
	\toprule
	\textbf{Policy}   & \textbf{Mechanism}                      & \textbf{Loses}             & \textbf{Used by}  \\
	\midrule
	Sliding window    & Remove oldest $n$ tokens                & Early conversation context & Most chatbots     \\
	Token budget      & Assign each source a limit              & Low-priority sources       & Custom agents     \\
	Relevance scoring & Score each message, evict lowest        & Context-dependent quality  & Research systems  \\
	Summarisation     & Replace $n$ old messages with a summary & Detail, nuance             & LangChain, MemGPT \\
	\bottomrule
\end{longtable}

\begin{principle}[Recency--Relevance Trade-off]
	For most agent tasks, the most recent $n$ tokens contain 80\% of the information needed. A sliding window that keeps the last 4K tokens is often sufficient, augmented with retrieved long-term memory for specific facts.
\end{principle}

\subsection{Summarisation}

When early messages are evicted, their content is condensed:

\begin{lstlisting}[caption=Running summarisation for context management]
  # After every K turns, or when approaching token limit:
  summary_prompt = f"""Summarise the conversation so far.
  Include: user's goals, facts learned, decisions made,
  pending actions, and preferences.

  Conversation:
  {recent_messages}

  Previous summary: {existing_summary}
  """

  new_summary = llm(summary_prompt)
  messages = [
    system_prompt,
    {"role": "system", "content": f"[Summary: {new_summary}]"},
    *recent_messages[-K:]     # Keep only last K turns
  ]
\end{lstlisting}

\subsection{Structured Working Memory}

An alternative to unstructured summarisation: extract key information into a structured schema that persists across turns.

\begin{lstlisting}[caption=Structured memory schema]
  AgentMemory ::= {
    goal:          string,
    completed:     Step[],
    pending:       Step[],
    facts_known:   Fact[],
    preferences:   { key: value },
    tool_states:   { tool_name: state },
    errors:        Error[],
    confidence:    number    // 0..1
  }
\end{lstlisting}

The structured memory is updated after every step and serialised into the prompt. This gives the \llm\ clean access to the full state without parsing the conversation history.

\section{Long-Term Memory}

Information that outlives a single session must be stored externally.

\subsection{Vector Memory}

The most common long-term memory architecture uses embedding-based retrieval:

\begin{lstlisting}[caption=Vector memory store]
  class VectorMemory:
      def __init__(self):
          self.entries = []    // [{text, embedding, metadata}]

      def add(self, text: str, metadata: dict):
          embedding = embed(text)
          self.entries.append({
              text: text,
              embedding: embedding,
              metadata: metadata,
              timestamp: now()
          })

      def query(self, query: str, top_k: int = 5):
          query_emb = embed(query)
          scores = cosine_similarity(query_emb,
                                     [e.embedding for e in self.entries])
          top_indices = argsort(scores)[-top_k:]
          return [self.entries[i].text for i in top_indices]
\end{lstlisting}

\subsection{Memory Structure}

What should be stored in long-term memory?

\begin{itemize}
	\item \textbf{Conversation summaries:} condensed records of past sessions.
	\item \textbf{User facts:} preferences, personal information, past decisions.
	\item \textbf{Failed attempts:} what didn't work and why (reflection memory).
	\item \textbf{Successful patterns:} what worked and can be reused.
	\item \textbf{External knowledge:} documents, notes, reference material.
\end{itemize}

\subsection{Recall Strategies}

Retrieval is not a single query. Multiple strategies exist:

\begin{longtable}{p{3.5cm}p{4cm}p{4cm}p{3cm}}
	\toprule
	\textbf{Strategy} & \textbf{How it works}                        & \textbf{Best for}         & \textbf{Caveat}              \\
	\midrule
	Single query      & Embed current query, retrieve top-$k$        & Simple fact lookup        & Misses multi-faceted queries \\
	Multi-query       & Generate $n$ query variations, merge results & Broad exploration         & $n \times$ cost              \\
	Stepwise          & Retrieve at each reasoning step              & Dynamic information needs & Cumulative latency           \\
	Automatic         & \llm\ decides when to retrieve               & Adaptive context          & Unpredictable                \\
	\bottomrule
\end{longtable}

\subsection{Forgetting and Consolidation}

Long-term memory must be curated:

\begin{itemize}
	\item \textbf{Decay:} old entries are demoted in retrieval scores unless accessed.
	\item \textbf{Merge:} duplicate or contradictory facts are resolved.
	\item \textbf{Summarisation:} many related entries are condensed into one.
	\item \textbf{Archival:} entries older than a threshold are moved to cold storage.
\end{itemize}

\section{Episodic Memory for Agents}

Generative Agents~\cite{park2023generative} introduced the most complete episodic memory architecture:

\begin{lstlisting}[caption=Generative Agents memory stream]
  Memory Stream:
    Each agent maintains a list of observations (events, actions, reflections).
    Each observation has:
    - content: natural language description
    - timestamp: when it happened
    - importance: 1--10 (how noteworthy)
    - type: "event" | "action" | "reflection" | "observation"

  Retrieval:
    score(memory, query) =
      recency(memory) * w_r +    // More recent = higher score
      relevance(memory, query) * w_v +  // Embedding similarity
      importance(memory) * w_i    // More important = higher score

  Reflection:
    When importance scores surpass a threshold,
    the agent generates a high-level reflection
    that summarises patterns across memories.

  Retrieval results shape the agent's current plan and behaviour.
\end{lstlisting}

\section{Procedural Memory}

Procedural memory captures \emph{how} to do things:

\begin{itemize}
	\item \textbf{Tool recipes:} sequences of tool calls that accomplish common tasks.
	\item \textbf{Prompt templates:} reusable patterns for common reasoning steps.
	\item \textbf{Error recovery plans:} what to do when specific failures occur.
	\item \textbf{Decision trees:} when to use which strategy.
\end{itemize}

Procedural memory is often encoded as few-shot examples in the system prompt or as a separate \llm\ call that retrieves the relevant procedure before the main reasoning step.

\section{Exercises}

\begin{exercise}
	Design a memory system for a personal assistant agent that manages calendar, email, and todo lists. Specify: what goes in working memory, what goes in long-term memory, and the retrieval strategy.
\end{exercise}

\begin{exercise}
	Implement the three-factor scoring function from Generative Agents (recency $\times$ relevance $\times$ importance). Compare performance against pure cosine similarity on a test set of 100 memories.
\end{exercise}

\begin{exercise}
	A summarisation-based working memory system loses detail over time. Design an experiment to measure information loss as a function of conversation length. What metric would you use?
\end{exercise}

\begin{exercise}
	Compare structured working memory (a schema) vs.\ unstructured (running summary). Which is better for: (a) a code generation agent, (b) a customer support agent, (c) a creative writing assistant?
\end{exercise}
